\documentclass[11pt,a4paper]{article}

\usepackage[utf8]{inputenc}
\usepackage[T1]{fontenc}
\usepackage{geometry}
\geometry{margin=1in, top=1.2in, bottom=1.2in}
\usepackage{hyperref}
\usepackage{booktabs}
\usepackage{tabularx}
\usepackage{graphicx}
\usepackage{listings}
\usepackage{xcolor}
\usepackage{enumitem}

\definecolor{codebackground}{rgb}{0.96,0.96,0.96}
\definecolor{keywordcolor}{rgb}{0.13,0.13,0.5}
\definecolor{stringcolor}{rgb}{0.6,0.1,0.1}

\lstset{
    backgroundcolor=\color{codebackground},
    basicstyle=\ttfamily\small,
    breaklines=true,
    captionpos=b,
    commentstyle=\color{gray},
    keywordstyle=\color{keywordcolor}\bfseries,
    stringstyle=\color{stringcolor},
    frame=lines,
    showstringspaces=false
}

\title{\textbf{Aroha: The Network for Autonomous Agent Protocol}\\ \Large{Whitepaper 1.0 (Draft)}}
\author{Aroha Working Group}
\date{April 28, 2026}

\begin{document}

\maketitle
\tableofcontents
\newpage

\section{Executive Summary}
As the world transitions to an AI-driven economy, autonomous agents are being built and deployed at an unprecedented scale. However, today's multi-agent systems face a structural limitation: \textbf{every agent is an island}. An agent built by one enterprise cannot natively discover, authenticate, or transact with an agent from another vendor. Today's orchestration frameworks (like LangGraph or CrewAI) are designed for local, in-framework coordination. Protocols like MCP connect agents to data and tools but do not facilitate secure agent-to-agent commerce.

The \textbf{Aroha Protocol (Aroha)} solves this. Aroha is an open, decentralized, vendor-neutral protocol stack built to provide \textbf{discovery, verifiable identity, and atomic transactions} for AI agents. It operates similarly to how TCP/IP and HTTP layered the web, enabling secure, machine-to-machine interactions globally without reliance on centralized API gateways.

Aroha bridges the gap natively with:
\begin{enumerate}[label=\textbf{\arabic*.}]
    \item \textbf{Cryptographic Identity:} \texttt{did:aroha:} via W3C Decentralised Identifiers and Ed25519.
    \item \textbf{Capability Discoverability:} Pluggable registries (EVM smart contracts + IPFS manifests).
    \item \textbf{Multi-Agent Atomic Transactions:} The Saga pattern (\texttt{Reserve $\rightarrow$ Commit $\rightarrow$ Cancel}).
    \item \textbf{Agent Commerce:} A pluggable settlement, RBAC, and escrow layer for B2B/B2C exchange.
\end{enumerate}

\section{The Problem Statement}
Current AI infrastructure lacks standard communication primitives for multi-agent capabilities. Building cross-vendor agent systems today results in three catastrophic points of failure:

\begin{table}[h]
    \centering
    \renewcommand{\arraystretch}{1.5}
    \begin{tabularx}{\textwidth}{@{} >{\bfseries}p{3cm} p{5.5cm} X @{}}
        \toprule
        Failure Dimension & Phenomenon & Systemic Impact \\
        \midrule
        1. Discovery & No standard protocol for publishing or discovering agent capabilities. Every connection must be hardcoded by a human. & Brittle, unscalable point-to-point integrations. Agents cannot organically find partners for tasks without manual human intervention. \\
        2. Verification & No cryptographic proof of identity. Standard APIs trust implicit callers or rely on easily stolen static API keys. & Rampant vulnerability to impersonation, man-in-the-middle, and replay attacks. Enterprises refuse to let autonomous agents communicate externally. \\
        3. Transactions & No capability safely runs multi-step actions across different vendor systems. If step 3 of a workflow fails, steps 1 and 2 remain finalized. & Catastrophic failure states with duplicate charges, partial database execution, and no global rollback or compensation mechanisms. \\
        \bottomrule
    \end{tabularx}
    \caption{The Three Dimensions of Multi-Agent Infrastructure Failure}
\end{table}

\section{The Aroha Solution \& Architecture}
Aroha utilizes a \textbf{six-layer kernel} combined with an opt-in extension system for commerce and settlement. The kernel strictly isolates routing and trust from business logic.

\subsection{Architecture Stack Diagram}
\begin{verbatim}
┌────────────────────────────────────────────────────────────────────┐
│  PROTOCOL EXTENSIONS (opt-in, never statically imported by kernel) │
│                                                                    │
│  @aroha-sdk/credentials   RBAC middleware, human credential issuance    │
│  @aroha-sdk/settlement    Pluggable billing: NullSettlement → Escrow    │
│  @aroha-sdk/commerce      B2B · B2C · C2B commerce, audit, consent      │
└────────────────────────────────────────────────────────────────────┘
                              ↑ mount via middleware / inject
┌───────────────────────────────────────────────┐
│  Layer 5 — Preference Context                 │ (budget, consent, signs)
├───────────────────────────────────────────────┤
│  Layer 4 — Orchestration                      │ (reserve/commit saga)
├───────────────────────────────────────────────┤
│  Layer 3 — Messaging                          │ (signed JSON-LD envelopes)
├───────────────────────────────────────────────┤
│  Layer 2 — Registry                           │ (EVM + IPFS manifests)
├───────────────────────────────────────────────┤
│  Layer 1 — Identity & Auth                    │ (did:aroha, Ed25519)
├───────────────────────────────────────────────┤
│  Layer 0 — Transport                          │ (HTTPS, wss://, TLS 1.3)
└───────────────────────────────────────────────┘
\end{verbatim}

\subsection{Detailed Protocol Layers}

The separation of layers ensures developers can mix and match components depending on their trust requirements (e.g., swapping public EVM discovery for private HTTP registries).

\begin{description}
    \item[Layer 0 (Transport)] Defines standard HTTP interactions. Agents expose \texttt{GET /.well-known/aroha-agent.json} for identity discovery, \texttt{POST /aroha/v1} for stateless JSON-LD message reception, and \texttt{wss /aroha/stream} for asynchronous progress streaming.
    \item[Layer 1 (Identity)] Every agent generates an Ed25519 keypair and creates a \texttt{did:aroha:<unique-id>} identifier. The DID acts as the root of trust, making spoofing mathematically impossible.
    \item[Layer 2 (Registry)] Instead of hardcoded URLs, capabilities are published as JSON-LD manifests to IPFS. The resulting CID is anchored to a registry (either managed HTTP or EVM-compatible Smart Contracts). This enables semantic querying: "Find an agent capable of \texttt{book-flight}."
    \item[Layer 3 (Messaging)] Transmitted payloads take the shape of canonical JSON-LD envelopes signed securely. Aroha strictly limits interactions to 15 core message types bridging queries, logic, and commercial negotiation.
    \item[Layer 4 (Orchestration)] Integrates the Saga transaction pattern as a protocol primitive.
    \item[Layer 5 (Preference Context)] A native context payload encompassing user constraints such as budgets, limits, and explicit tokenized consent rules.
\end{description}

\section{Key Protocol Features}

\subsection{Machine-to-Machine Messaging Structure}
Aroha envelopes apply zero-trust principles. Every payload contains an Ed25519 signature over the canonicalized subset of the message. 

\begin{lstlisting}[caption=Example Signed Aroha Envelope, basicstyle=\ttfamily\small]
{
  "@context": [
    "https://aroha-labs.com/contexts/v1",
    "https://w3id.org/security/suites/ed25519-2020/v1"
  ],
  "id": "urn:uuid:7f6a...",
  "type": "ArohaRequest",
  "from": "did:aroha:personal-agent",
  "to": "did:aroha:flight-provider",
  "created": "2026-04-28T08:00:00Z",
  "expires": "2026-04-28T08:05:00Z",
  "nonce": "e3b0c44298fc1c149afbf4c8996fb924",
  "body": {
    "capability": "book-flight",
    "params": { "destination": "NRT", "date": "2026-05-10" }
  },
  "proof": {
    "type": "Ed25519Signature2020",
    "proofValue": "<base64url-signature>"
  }
}
\end{lstlisting}

\subsection{Cross-Agent Sagas (Atomic Rollbacks)}
A primary innovation inside Aroha is atomic cross-agent state changes. Workflows adhere to a strict three-phase commit pattern: \texttt{Reserve $\rightarrow$ Commit $\rightarrow$ Cancel}. LIFO (Last-In, First-Out) compensation guarantees no disjointed databases or stranded funds.

\begin{verbatim}
[Saga State Machine - Multi-Agent Rollback]

             [ ArohaReserve Sent ]
[ IDLE ] ---------------------------> [ RESERVING ]
                                            │
                                            ├─ (Success / Ack) ─> [ RESERVED ]
                                            │                            │
                                            ├─ (Failure / T.O) ─> [ FAILED ]
                                            │                            │
           [ COMMIT INITIATED ]             │                            │
           [ RESERVED ] --------------------+        [ REVERT INITIATED ]│
                │                                    [ RESERVED/COMMITTING ]
                ├─ (Success) ──> [ COMMITTED ]               │
                │                                            ▼
                ├─ (Failure) ──> [ COMPENSATING ] ──> [ COMPENSATED (Rollback) ]
\end{verbatim}

\subsection{Invisible Web3 Settlement \& Commerce}
Agent economy requires pure M2M (machine-to-machine) micro-payments. Aroha isolates billing logic via a pluggable abstraction known as the \texttt{WalletService}. End-users never handle wallets, gas, or signatures directly.

\begin{itemize}
    \item \textbf{Null / Free Tier:} Direct trust networks bypass settlement modules entirely.
    \item \textbf{Invisible Pre-Pay:} Orchestrators generate ephemeral session keys logic utilizing paymasters to sponsor chain fees. 
    \item \textbf{On-Chain Escrow:} Tokenized value is escrowed pending a valid, cryptographic cryptographic receipt that unlatches funds only on \texttt{ArohaCommitAck}.
\end{itemize}

\section{Security \& Threat Model}
Deploying fully autonomous agent systems mandates a Zero-Trust security approach enforced at the Transport boundaries (Layer 0 \& Layer 1) before business logic fires.

\begin{table}[h]
    \centering
    \renewcommand{\arraystretch}{1.4}
    \begin{tabularx}{\textwidth}{@{} l X @{}}
        \toprule
        \textbf{Threat Vector} & \textbf{Aroha Mitigation Strategy} \\
        \midrule
        \textbf{Impersonation} & Every message requires an Ed25519 cryptographic signature. \\
        \textbf{Replay Attacks} & Per-message 32-byte nonces stored in a rolling 5-minute TTL cache matched rigorously against the \texttt{expires} field. \\
        \textbf{Payload Tampering} & Signature verifies a canonicalized structure containing sorted payload keys, rejecting any mismatched hashing instantly. \\
        \textbf{Eavesdropping/MITM} & All Layer 0 transit defaults to TLS 1.3+. \\
        \textbf{Privilege Escalation} & RBAC policy definitions evaluate the sender URI mapping \texttt{CredentialTokens} sequentially prior to calling the core handlers. \\
        \bottomrule
    \end{tabularx}
    \caption{Zero-Trust Threat Mitigations}
\end{table}

\section{Trust Levels and Reputation}
Cryptographic identity alone ensures verification, not necessarily performance or commercial reliability. Aroha evaluates capability and past behavior via progressive, oracle-driven trust levels:

\begin{table}[h]
    \centering
    \renewcommand{\arraystretch}{1.3}
    \begin{tabularx}{\textwidth}{@{} >{\bfseries}p{1.8cm} p{6cm} X @{}}
        \toprule
        Level & Protocol Requirement & Allowed Actor Operations \\
        \midrule
        Lvl 0 & Unverified / Anonymous & Read Only, IPFS Registry Query. \\
        Lvl 1 & On-chain registration + Stake limits & Agent Discovery and price Negotiation logic. \\
        Lvl 2 & W3C Verifiable Credential Identity & Low-value \texttt{Reserve} allocations. \\
        Lvl 3 & On-Chain Reputation > 80\% & High-value financial commits, large tokenized logic operations. \\
        Lvl 4 & Explicit Authorization via User Signature & Administrative or full capability delegation functions. \\
        \bottomrule
    \end{tabularx}
    \caption{Agent Capability Trust Benchmarks}
\end{table}

\section{Summary}
Aroha standardizes interactions that transition autonomous agents from isolated, hardcoded logic bubbles to an expansive, decentralized, interoperable network. By combining W3C DIDs, JSON-LD message semantics, atomic Sagas, and native payment rails, Aroha serves as the bedrock protocol for the automation economy.

\end{document}